\section{Introduction}\label{sec:intro}

Six Birds Theory is a program that treats a theory as a fixed point of
an idempotent closure operator acting on descriptions of a system, and
studies how such closed descriptions arise
\cite{Tsiokos2026EmergenceCalculus}.  Its companion papers apply this
vocabulary to open problems (a glossary is in
\Cref{sec:glossary}): a \emph{conditional closure} of
a problem is a theorem deriving the problem's standard statement from
explicitly named premises.  Three companion papers state conditional
closures of three Clay Millennium problems in this sense
\cite{Fefferman2000NavierStokes,Bombieri2000Riemann,Cook2000PvsNP,CarlsonJaffeWiles2006MillenniumPrize}.
The Navier--Stokes paper \cite{Tsiokos2026NavierStokes} derives a
global solution satisfying the equations pointwise from a spectral
closure hypothesis and hypotheses on the initial data.  The Riemann-hypothesis paper
\cite{Tsiokos2026RH} derives the Riemann hypothesis from a supplied
Selberg recognition source and an identification with the classical
zeros.  The P-versus-NP paper \cite{Tsiokos2026PvNP} proves a fixed
machine-level translation between deciding SAT and computing a tagged
witness selector, and derives \(P\neq NP\) from an accepted negative
source.  In each case the premises are open.  None of the three is an
unconditional result, and this paper does not change that.  The three
are also not of one type: the Navier--Stokes closure is derived from its
hypotheses through problem-specific analysis, while the other two name
a load-bearing source as a premise.

The three arguments nevertheless have a common shape, and this paper
asks why.  It gives two kinds of answer.  The first is concrete: one
family of Gaussian observations can be read on all three mathematical
objects (\Cref{sec:gaussian-bridge,sec:triad-bridge}).  The second is
structural: a description, in the vocabulary of the framework, of how
proof attempts aimed at a target can stall, of how a closure can instead
name its load-bearing source, and of the detectors by which obstructions
are measured (\Cref{sec:foreclosure,sec:recognition,sec:calibration}).

\paragraph{The concrete construction.}  For a width \(a>0\) and a
complex centre \(z\), a complex Gaussian window on the Fourier side of
\(\mathbb R^{3}\) defines a linear observation of Fourier states.
Evaluated against the heat multiplier, it gives
\((\pi/(a+b))^{3/2}e^{-\kappa z^{2}}\) with \(b=4\pi^{2}\nu t\) and
\(\kappa=ab/(a+b)\).  Placing \(z\) at the spectral coordinate of a zero
\(\rho\) of the completed zeta function, the logarithmic excess of this
value is \(\kappa(\operatorname{Re}\rho-\frac12)^{2}\); summed with
reflection-invariant nonnegative integer weights over a finite window closed under
reflection in the critical line, it is \(\kappa\) times
the residual charge of the window in the sense of the
adequacy-residual companion \cite{Tsiokos2026XiCompanion}
(\Cref{thm:joint-gaussian-return}).  Vanishing of this excess on all
windows is a reformulation of the Riemann hypothesis, not a route to
it.  Applied to a mild solution of the Navier--Stokes equations the
same observation restates the mild equation, with the nonlinear Duhamel
term kept and not estimated.  On the assignment space of a CNF formula
the same scalar kernel gives a Gaussian mass that is positive exactly
for satisfiable formulas, with exact laws for adding a clause and for
fixing a variable (\Cref{thm:concrete-triad-gaussian-xi}); positivity
here holds for any \(\{0,1\}\)-valued filter and is not specific to
SAT.  On a reflected
zero window the constant probe reads zero; on a satisfiable formula it
reads the number of satisfying assignments.  These identities are
elementary, and they are formally verified in Lean
(\Cref{app:formalization}).  What the three settings share is a
formula, not a mechanism: the identities transfer no premise from one
problem to another, and no efficient evaluation of the Gaussian mass is
proved.

\paragraph{The structural account.}  The account has three parts.
Part~I (\Cref{sec:foreclosure}) concerns \emph{carriers}: a carrier is a
presentation through which a proof attempt tries to reach a target
\(X\), together with the obstruction it leaves, its \emph{residual}.
Four outcomes are distinguished: the residual closure points to a
proposition other than \(X\); it is equivalent to \(X\) but the attempt
can certify \(X\) only by auditing itself; it needs a bridge to \(X\)
that is unavailable; or its reduction ends at a family of matrix
elements whose decision it does not supply.  That these four outcomes
are exhaustive is the Coverage schema, which is not proved; the
theorems of Part~I derive consequences from explicitly assumed schemas
and are short logical arguments.

Part~II (\Cref{sec:recognition}) concerns closures that do not derive
their load-bearing content but name it as a premise, called
\emph{recognition-mode landings}.  It states, as four conditional
schemas, a seven-stage template for such closures, a pattern in the
outcome of searching for an internal derivation, the thesis that the
load-bearing content is named rather than derived, and a combination of
this with the foreclosure statement of Part~I.  The Riemann-hypothesis
and P-versus-NP applications illustrate the schemas; they do not prove
them.

Part~III (\Cref{sec:calibration}) concerns \emph{calibrations}, typed
detectors of obstructions such as low-degree polynomials, mutual
information or Nullstellensatz certificates.  It records a catalogue,
introduced here, of thirteen detector classes considered for the
P-versus-NP problem, a criterion for
when a new calibration is a relabelling of an old one, the direction of
bridges between calibrations, a table of which regions the catalogue
covers, and three named transports---\emph{functorial gates}---whose
proofs are open.  The results of Part~III are elementary facts about
these finite records; they say nothing about the performance of a
detector or about inclusions between complexity classes.

\paragraph{Status of the results.}  The paper distinguishes three
kinds of numbered statement (\Cref{sec:apparatus:opacity-guards}).
Theorems, propositions and corollaries are proved from the hypotheses
they state; in Parts~I and~III most of them are short logical arguments
or checks of declared finite data, and we say so where they occur.  Conditional schemas are not proved; they are displayed so
that other statements can assume them.  Some predicates, such as being
in framework scope or being load-bearing, are primitive: they are named
and not defined here, and results about them establish only their
logical shape.  \Cref{app:formalization} records, for every numbered
statement, its formal counterpart in Lean.  The only computations in
the paper are in \Cref{sec:gaussian-bridge,sec:triad-bridge}, and they
are elementary; their contribution is that the definitions are aligned
across the three settings and formally verified.  The rest of the paper
is a vocabulary for comparing conditional closures, and a reader may
take it as such.
A glossary of the framework's terms is in
\Cref{sec:glossary}.  The companion papers are preprints.

\paragraph{Relation to the series.}  The paper uses the closure
calculus of \FOne\ \cite{Tsiokos2026EmergenceCalculus}, the
admissibility setting of \FTwo\ \cite{Tsiokos2026FoundationsII}, and the
audit calculus of \FThree\ \cite{Tsiokos2026FoundationsIII}.  The three
applications \cite{Tsiokos2026NavierStokes,Tsiokos2026RH,Tsiokos2026PvNP}
are cited as they stand; the recognition sources of the two
trace-state-only applications are named in those applications.  The
self-dual trace confinement companion \cite{Tsiokos2026SDTC} supplies
the confinement squeeze and its setting, the hiddenness companion
\cite{Tsiokos2026CSL} the placements and exposure classifications used
here, the adequacy-residual companion \cite{Tsiokos2026XiCompanion} supplies
the residual, and the emergence companion \cite{Tsiokos2026NeedleKiller}
supplies the distinction between native and layer-dissolving probes.
This paper is downstream of all of these and replaces none of their
arguments.

\paragraph{Scope.}  The Birch--Swinnerton-Dyer conjecture
\cite{Wiles2000BSD} is outside the scope of the paper, and the Hodge
conjecture \cite{Deligne2000Hodge} appears only as a prediction.  Nothing
here asserts an unconditional proof of any Millennium problem
(\Cref{sec:scope}).

\paragraph{Organization.}  \Cref{sec:apparatus} fixes the vocabulary.
\Cref{sec:foreclosure,sec:recognition} contain Parts~I and~II, and
\Cref{sec:validation} places the three applications in that scheme.
\Cref{sec:gaussian-bridge,sec:triad-bridge} contain the Gaussian
construction, and \Cref{sec:calibration} Part~III.
\Cref{sec:audit} checks the description of the companion papers,
\Cref{sec:predictions} lists predictions and open problems, and
\Cref{sec:scope} collects the limitations.
\Cref{app:formalization} describes the formal verification and
\Cref{app:version-history} the changes between versions.
